% GENERATED FILE. Edit canonical lessons, then run npm run generate:books.
% canonical-source-hash: fnv1a64:247fe5302d037a43
% canonical-lessons: RU-C63-white, RU-C63-black, RU-C63-red, RU-C63-blue, RU-C63-green

\chapter{White, Black, Red, Blue, Green}
\label{ch:ru-white-black-red-blue-green}
\hlchaptermodality{63}

\begin{chapteropening}
\textbf{By the end of this chapter:} \emph{I can say white, black, red, blue and green.}

Complete the last lesson of chapter 63: I can say white, black, red, blue and green.
\end{chapteropening}

\section[\texorpdfstring{\ru{белый} (bélyy)}{белый (bélyy)}]{\ru{белый} (bélyy) — white}

\label{lesson:RU-C63-white}



\begin{quote}
Before the new one: say the Russian for a duck, then the Russian for a goat.
\end{quote}

\subsection*{You'll want to know: \ru{белый}}
\textbf{\ru{белый}} (\emph{bélyy}) — \textquotedblleft{}white\textquotedblright{}.

\begin{cousinweb}
\textbf{\ru{белый}} $\leftarrow$ Proto-Slavic \emph{*bělŭ} — inherited. Vasmer relates it to Lithuanian \emph{bãlas}, \textquotedblleft{}white\textquotedblright{}, Latvian \emph{bãls}, \textquotedblleft{}pale\textquotedblright{}, Greek \emph{phalós}, \textquotedblleft{}white\textquotedblright{}, and Sanskrit \emph{bhāti}, \textquotedblleft{}shines\textquotedblright{}, from PIE \emph{*b\textsuperscript{h}el-}, \textquotedblleft{}to shine\textquotedblright{}.
\end{cousinweb}

\begin{grammarlens}[title={What you've built}]
One more word for this chapter.
\end{grammarlens}

\subsection*{Your turn}
\begin{itemize}
\raggedright
  \item \emph{Say it:} \emph{bélyy}
  \item \emph{Say it:} \emph{bélyy}, once more
  \item \emph{Say it:} say \emph{bélyy}
  \item \emph{Recall:} say the Russian for a duck, then the Russian for a goat, then say \emph{bélyy} again
\end{itemize}

\begin{tcolorbox}[breakable,colback=teal!4,colframe=teal!35!black,title={Before you move on}]
What does \ru{белый} mean? (\textquotedblleft{}White\textquotedblright{}.) Say it once more.
\end{tcolorbox}
\section[\texorpdfstring{\ru{чёрный} (chyórnyy)}{чёрный (chyórnyy)}]{\ru{чёрный} (chyórnyy) — black}

\label{lesson:RU-C63-black}



\begin{quote}
Before the new one: say the Russian for a goat, then the Russian for white.
\end{quote}

\subsection*{You'll want to know: \ru{чёрный}}
\textbf{\ru{чёрный}} (\emph{chyórnyy}) — \textquotedblleft{}black\textquotedblright{}.

\begin{cousinweb}
\textbf{\ru{чёрный}} $\leftarrow$ Proto-Slavic \emph{*čĭrnŭ}, from \emph{*čĭrxnŭ} — inherited. Vasmer relates it to Old Prussian \emph{kirsnan}, \textquotedblleft{}black\textquotedblright{}, and Sanskrit \emph{kṛṣṇá-}, \textquotedblleft{}black\textquotedblright{}, the name of Krishna.
\end{cousinweb}

\begin{grammarlens}[title={What you've built}]
One more word for this chapter.
\end{grammarlens}

\subsection*{Your turn}
\begin{itemize}
\raggedright
  \item \emph{Say it:} \emph{chyórnyy}
  \item \emph{Say it:} \emph{chyórnyy}, once more
  \item \emph{Say it:} say \emph{chyórnyy}
  \item \emph{Recall:} say the Russian for a goat, then the Russian for white, then say \emph{chyórnyy} again
\end{itemize}

\begin{tcolorbox}[breakable,colback=teal!4,colframe=teal!35!black,title={Before you move on}]
What does \ru{чёрный} mean? (\textquotedblleft{}Black\textquotedblright{}.) Say it once more.
\end{tcolorbox}
\section[\texorpdfstring{\ru{красный} (krásnyy)}{красный (krásnyy)}]{\ru{красный} (krásnyy) — red}

\label{lesson:RU-C63-red}



\begin{quote}
Before the new one: say the Russian for white, then the Russian for black.
\end{quote}

\subsection*{You'll want to know: \ru{красный}}
\textbf{\ru{красный}} (\emph{krásnyy}) — \textquotedblleft{}red\textquotedblright{}.

\begin{cousinweb}
\textbf{\ru{красный}} is built on \emph{\ru{краса}}, \textquotedblleft{}beauty\textquotedblright{}, and first meant \textquotedblleft{}beautiful\textquotedblright{}, as in \emph{\ru{прекрасный}}, \textquotedblleft{}lovely\textquotedblright{}; the name of Moscow's Red Square first meant \textquotedblleft{}the beautiful square\textquotedblright{}. Vasmer thinks \emph{\ru{краса}} is probably related to Old Norse \emph{hrósa}, \textquotedblleft{}to boast\textquotedblright{}.
\end{cousinweb}

\begin{grammarlens}[title={What you've built}]
One more word for this chapter.
\end{grammarlens}

\subsection*{Your turn}
\begin{itemize}
\raggedright
  \item \emph{Say it:} \emph{krásnyy}
  \item \emph{Say it:} \emph{krásnyy}, once more
  \item \emph{Say it:} say \emph{krásnyy}
  \item \emph{Recall:} say the Russian for white, then the Russian for black, then say \emph{krásnyy} again
\end{itemize}

\begin{tcolorbox}[breakable,colback=teal!4,colframe=teal!35!black,title={Before you move on}]
What does \ru{красный} mean? (\textquotedblleft{}Red\textquotedblright{}.) Say it once more.
\end{tcolorbox}
\section[\texorpdfstring{\ru{синий} (síniy)}{синий (síniy)}]{\ru{синий} (síniy) — blue (dark blue)}

\label{lesson:RU-C63-blue}



\begin{quote}
Before the new one: say the Russian for black, then the Russian for red.
\end{quote}

\subsection*{You'll want to know: \ru{синий}}
\textbf{\ru{синий}} (\emph{síniy}) — \textquotedblleft{}blue (dark blue)\textquotedblright{}.

\begin{cousinweb}
\textbf{\ru{синий}} $\leftarrow$ Proto-Slavic \emph{*sinĭ} — inherited. Vasmer relates it to \emph{\ru{сиять}}, \textquotedblleft{}to shine\textquotedblright{}, or to \emph{\ru{сивый}}, \textquotedblleft{}grey\textquotedblright{}, and to Sanskrit \emph{śyāmá-}, \textquotedblleft{}dark\textquotedblright{}, and Lithuanian \emph{šývas}, \textquotedblleft{}light-coloured\textquotedblright{}.
\end{cousinweb}

\begin{grammarlens}[title={What you've built}]
One more word for this chapter.
\end{grammarlens}

\subsection*{Your turn}
\begin{itemize}
\raggedright
  \item \emph{Say it:} \emph{síniy}
  \item \emph{Say it:} \emph{síniy}, once more
  \item \emph{Say it:} say \emph{síniy}
  \item \emph{Recall:} say the Russian for black, then the Russian for red, then say \emph{síniy} again
\end{itemize}

\begin{tcolorbox}[breakable,colback=teal!4,colframe=teal!35!black,title={Before you move on}]
What does \ru{синий} mean? (\textquotedblleft{}Blue (dark blue)\textquotedblright{}.) Say it once more.
\end{tcolorbox}
\section[\texorpdfstring{\ru{зелёный} (zelyónyy)}{зелёный (zelyónyy)}]{\ru{зелёный} (zelyónyy) — green}

\label{lesson:RU-C63-green}



\begin{quote}
Before the new one: say the Russian for red, then the Russian for blue.
\end{quote}

\subsection*{You'll want to know: \ru{зелёный}}
\textbf{\ru{зелёный}} (\emph{zelyónyy}) — \textquotedblleft{}green\textquotedblright{}.

\begin{cousinweb}
\textbf{\ru{зелёный}} $\leftarrow$ Proto-Slavic \emph{*zelenŭ} — inherited. Vasmer relates it to Lithuanian \emph{žãlias}, \textquotedblleft{}green\textquotedblright{}, and \emph{žélti}, \textquotedblleft{}to grow over\textquotedblright{}, and further to \emph{\ru{золото}}, \textquotedblleft{}gold\textquotedblright{}, and Sanskrit \emph{hári-}, \textquotedblleft{}yellow\textquotedblright{}. German \emph{gelb}, \textquotedblleft{}yellow\textquotedblright{}, comes from the same PIE root, \emph{*ǵ\textsuperscript{h}elh\textsubscript{3}-}, \textquotedblleft{}to be green, yellow\textquotedblright{}.
\end{cousinweb}

\begin{grammarlens}[title={What you've built}]
One more word for this chapter.
\end{grammarlens}

\subsection*{Your turn}
\begin{itemize}
\raggedright
  \item \emph{Say it:} \emph{zelyónyy}
  \item \emph{Say it:} \emph{zelyónyy}, once more
  \item \emph{Say it:} say \emph{zelyónyy}
  \item \emph{Recall:} say the Russian for red, then the Russian for blue, then say \emph{zelyónyy} again
\end{itemize}

\begin{tcolorbox}[breakable,colback=teal!4,colframe=teal!35!black,title={Before you move on}]
What does \ru{зелёный} mean? (\textquotedblleft{}Green\textquotedblright{}.) Say it once more.
\end{tcolorbox}
